\section{Fun with Propositional Logic}

A propositional well-formed formula (WFF) is defined inductively over a signature of variable symbols $\mathcal{V}$ and logical connectives $\{\neg, \wedge, \vee, \to\}$:
\begin{itemize}
    \item Every propositional variable $p \in \mathcal{V}$ is a formula.
    \item If $\varphi$ is a formula, then $\neg \varphi$ is a formula.
    \item If $\varphi$ and $\psi$ are formulas, then $(\varphi \wedge \psi)$, $(\varphi \vee \psi)$, and $(\varphi \to \psi)$ are formulas.
\end{itemize}

We want to write a predicate \texttt{isForm(L,F)} that checks if $F$ is a well-formed formula given the list of variables $L$.

\begin{codeblock}
isForm(L, var(X)) :- 
    string(X), 
    is_member(X, L).

isForm(L, neg(X)) :- 
    isForm(L, X).

isForm(L, disj(X, Y)) :- 
    isForm(L, X), 
    isForm(L, Y).

isForm(L, conj(X, Y)) :- 
    isForm(L, X), 
    isForm(L, Y).

isForm(L, impl(X, Y)) :- 
    isForm(L, X), 
    isForm(L, Y).
\end{codeblock}

Let's test this:

\begin{codeblock}
% Valid negation
?- isForm(["a", "b"], neg(var("a"))).
true.

% Variable 'c' is not in the declared alphabet ["a", "b"]
?- isForm(["a", "b"], neg(var("c"))).
false.

% Arity mismatch: impl requires 2 arguments
?- isForm(["a", "b"], conj(var("a"), impl(var("b")))).
false.

% Nested compound formula
?- isForm(["p", "q", "r"], impl(var("p"), disj(var("q"), var("r")))).
true.
\end{codeblock}

\minititle{Simplify}
In propositional logic, the main connectives are $\neg$ and $\land$. Every other connectives can be written using negation and conjunction.

\[
A \vee B \equiv \neg(\neg A \wedge \neg B)
\]

\[
A \Rightarrow B \equiv \neg A \vee B \equiv \neg(A \wedge \neg B)
\]


We want a predicate \texttt{simplify(Formula, Simplified)} that takes an arbitrary formula and transforms it into an equivalent formula containing only \texttt{var}, \texttt{neg}, and \texttt{conj}.


\begin{codeblock}
% Base Case: a variable is already simplified
simplify(var(X),var(X)).

% A negation of something is simplified if that something is simplified
simplify(neg(X), neg(Y)):-simplify(X,Y).

% A conjunction of two things is simplified if the two things are
simplify(conj(X,Y),conj(X1,Y1)):-simplify(X,X1), simplify(Y,Y1).

% Disjunction
simplify(disj(X,Y),neg(conj(neg(X1),neg(Y1)))):-
    simplify(X,X1),
    simplify(Y,Y1).

% Implication
simplify(impl(X,Y), neg(conj(X1, neg(Y1)))):-
    simplify(X,X1),
    simplify(Y,Y1).
\end{codeblock}

Let's test this:

\begin{codeblock}
?- simplify(disj(var("a"), var("b")), X).
X = neg(conj(neg(var("a")), neg(var("b")))).

?- simplify(impl(var("a"), var("b")), X).
X = neg(conj(var("a"), neg(var("b")))).
\end{codeblock}

\minititle{Evaluation}
Now we turn to \defemph{semantics}: computing the truth value of a formula under a given variable assignment (valuation) $\nu : \mathcal{V} \to \{0, 1\}$.

We represent an assignment as a list of key-value associations:
\[
[\texttt{assign("a", 1)},\, \texttt{assign("b", 0)}]
\]
We want a predicate \texttt{eval(Assignments, Formula, Value)} where $\text{Value} \in \{0, 1\}$.

\questionbox{
How do we evaluate logical operators cleanly in Prolog without writing multi-clause truth tables for every single connective?
}

\titledbox{Algebraic Encoding of Boolean Logic}{
\begin{align*}
\nu(\neg A) &= 1 - \nu(A) \\
\nu(A \wedge B) &= \nu(A) \times \nu(B) \\
\nu(A \vee B) &= \nu(A) + \nu(B) - (\nu(A) \times \nu(B)) \\
\nu(A \Rightarrow B) &= \nu(\neg A \vee B) = (1 - \nu(A)) + \nu(B) - ((1 - \nu(A)) \times \nu(B))
\end{align*}
}

\titledbox{Sanity Check - Disjunction}{
\begin{itemize}
    \item $0 \vee 0 = 0 + 0 - 0 = \mathbf{0}$
    \item $0 \vee 1 = 0 + 1 - 0 = \mathbf{1}$
    \item $1 \vee 0 = 1 + 0 - 0 = \mathbf{1}$
    \item $1 \vee 1 = 1 + 1 - 1 = \mathbf{1}$
\end{itemize}
}

\ideabox{
We use this algebraic encoding to evaluate formulas!
}


\begin{codeblock}
% Base Case (variable) -> just a lookup
eval(L, var(X), Y) :-
    is_member(assign(X, Y), L).

% Negation
eval(L, neg(X), Z) :-
    eval(L, X, V),
    Z is 1 - V.

% Conjunction
eval(L, conj(X, Y), V) :-
    eval(L, X, X1),
    eval(L, Y, Y1),
    V is X1 * Y1.

% Disjunction
eval(L, disj(X, Y), V) :-
    eval(L, X, X1),
    eval(L, Y, Y1),
    Plus is X1 + Y1,
    Mult is X1 * Y1,
    V is Plus - Mult.

% Implication
eval(L, impl(X, Y), V) :-
    eval(L, X, X1),
    eval(L, Y, Y1),
    NotX is 1 - X1,
    Plus is NotX + Y1,
    Mult is NotX * Y1,
    V is Plus - Mult.
\end{codeblock}

Let's test this:

\begin{codeblock}
?- eval([assign("a", 1), assign("b", 0)], conj(var("a"), var("b")), V).
V = 0.

?- eval([assign("p", 1), assign("q", 0)], impl(var("p"), var("q")), V).
V = 0.

?- eval([assign("p", 0), assign("q", 0)], impl(var("p"), var("q")), V).
V = 1.

?- eval([assign("a", 0), assign("b", 0), assign("c", 1)], 
        conj(neg(var("a")), disj(var("b"), var("c"))), V).
V = 1.
\end{codeblock}